\chapter{Unit 3: Fundamentals of Integral Calculus}

\section{Section 1: 20 Solved Comprehensive Subjective Problems}

\begin{subjectiveq}{3.1}
Evaluate the definite integral:
\begin{equation}
I = \int_0^\pi \frac{x \sin x}{1 + \cos^2 x}\,dx
\end{equation}
\end{subjectiveq}

\begin{solbox}
1. \textbf{Apply King's Property} $\int_0^a f(x)\,dx = \int_0^a f(a-x)\,dx$:
\begin{equation}
I = \int_0^\pi \frac{(\pi - x) \sin(\pi - x)}{1 + \cos^2(\pi - x)}\,dx = \int_0^\pi \frac{(\pi - x)\sin x}{1 + \cos^2 x}\,dx \quad \text{--- (2)}
\end{equation}
2. \textbf{Add (1) and (2)}:
\begin{align}
2I &= \int_0^\pi \frac{[x + (\pi - x)]\sin x}{1 + \cos^2 x}\,dx = \pi \int_0^\pi \frac{\sin x}{1 + \cos^2 x}\,dx
\end{align}
3. \textbf{Substitution}:
Let $u = \cos x \implies du = -\sin x\,dx$.
When $x = 0$, $u = 1$; when $x = \pi$, $u = -1$:
\begin{align}
2I &= \pi \int_1^{-1} \frac{-du}{1 + u^2} = \pi \int_{-1}^1 \frac{du}{1 + u^2} = \pi \Big[ \tan^{-1}u \Big]_{-1}^1 \\
&= \pi \left( \frac{\pi}{4} - \left(-\frac{\pi}{4}\right) \right) = \pi \left( \frac{\pi}{2} \right) = \frac{\pi^2}{2}
\end{align}
Therefore, $I = \frac{\pi^2}{4}$.
\end{solbox}

\section{Section 2: 50 Solved Multiple-Choice Questions (MCQs)}
\mcqitem{1}{The value of $\int_0^{\pi/2} \frac{\sin x}{\sin x + \cos x}\,dx$ is:}{$\pi$}{$\pi/2$}{$\pi/4$}{$0$}
\mcqitem{2}{If $f(x)$ is an odd function, then $\int_{-a}^a f(x)\,dx$ equals:}{$2\int_0^a f(x)dx$}{$0$}{$a$}{$-a$}
\mcqitem{3}{The value of $\int \ln x\,dx$ is:}{$1/x + C$}{$x\ln x - x + C$}{$x\ln x + x + C$}{$\ln x + C$}

\section{Section 3: Answer Key \& Technical Rationales}
\begin{table}[htbp]
\centering
\caption{\textbf{Unit 3: Answer Key \& Rationales}}
\vspace{2mm}
\begin{tabularx}{\textwidth}{l c X l c X}
\toprule
\textbf{Q\#} & \textbf{Ans} & \textbf{Rationale} & \textbf{Q\#} & \textbf{Ans} & \textbf{Rationale} \\
\midrule
1 & \textbf{(c)} & By King's property $2I = \int_0^{\pi/2} 1\,dx = \pi/2 \implies I = \pi/4$ & 3 & \textbf{(b)} & Integration by parts $\int 1 \cdot \ln x dx = x\ln x - x + C$ \\
2 & \textbf{(b)} & Odd symmetry over $[-a, a]$ integrates identically to zero & & & \\
\bottomrule
\end{tabularx}
\end{table}
